\section{还没做完的，以及接下来做什么}

\subsection{明确决定不做的}

下面四项不是“还差一点”，是本次决定不做，理由写在这里。

\textbf{语音。} 已经整体停用（\code{VOICE\_FEATURE=false}），控件隐藏，代码留在开关后面。原因是它连续三次判断错方向，而我没有办法独立验证音频——无头浏览器没有声卡，每一次判断都要占用使用者的时间。恢复语音之前这一块没有可观察对象，也无法验收。

\textbf{联网 PVP。} 本地对战是同一台设备上分屏轮流操作，服务端仍然接受客户端传来的状态快照，不是权威对局服务。真实竞技要求服务端下发可见状态和模式，否则改一个字段就能绕过限制。这条差距不是工程量问题，是架构问题。

\textbf{LLM 权重训练。} 本机是 arm64，没有可用的 CUDA 后端，也没有云上 GPU 预算。仓库里记录的 \code{cloud\_budget} 是 \code{not supplied}。

\textbf{训练期原始轨迹的完整归档。} 核对后确认核心要求已经满足：三个种子的独立测试结果一致，切分与轨迹摘要在。唯一缺的是训练期原始轨迹的完整落盘，属于归档事项。这个实验的说服力取决于它是什么（1152 个参数、英文二选一、冻结骨干），而不是证据的厚度，所以我决定不做。

\subsection{受条件限制，做不了的}

这些不是决定不做，是缺外部条件。

\begin{table}[h]
\centering
\small
\begin{tabular}{L{4.6cm}L{8.6cm}}
\toprule
项目 & 缺什么 \\
\midrule
无提示迁移验证 & 需要真人被试在多个情境里重复观察；机制已实现，缺的是数据 \\
人工试玩与学习效果观察 & 需要真人被试。现在没有真人学习效果的数据，模拟奖励不等于真人学习 \\
分屏双人完整验收 & 需要两名真人打多局；分屏逻辑我用脚本模拟点击验证过 \\
浏览器慢响应验收 & 需要可控的慢网环境。本机延迟只有 2--4 秒，不够 \\
独立盲标的检索评测集 & 现在 128 条是我自写自标，不是独立盲标，绝对数字偏乐观 \\
独立大样本模型质量评测 & 需要独立标注集与 API 预算；现有真实样本不构成质量结论 \\
\bottomrule
\end{tabular}
\end{table}

慢网那一条要补一句：使用者曾在公交车上用不稳定的连接真实游玩过，过期建议的场景是实际发生过的，不是假想；缺的只是能在浏览器里稳定复现它的环境。客户端一侧已经用替换请求的方式验证了 6 项。语音的听感验证也挂在这里——本机没有音频输出，我无法自验。

\subsection{接下来的动作}

按可核查的方式写，每项都带完成标志。

\begin{table}[h]
\centering
\small
\begin{tabular}{L{6.4cm}L{5.6cm}L{1.4cm}}
\toprule
动作 & 完成标志 & 预计 \\
\midrule
让生成脚本从表里推出那句话 & 改表之后不再需要手改文案 & 1 小时 \\
给提示正文加禁止词表与测试 & “启发式评分”“分支”不再出现在玩家读到的句子里 & 半天 \\
按新政策重新标注评测集并复测 & 重新标注在看见结果之前完成 & 半天 \\
给检索加最低分阈值与拒答再对照 & 负例拒答率显著高于 56.25\% & 半天 \\
补独立标注者，重做检索盲测 & 不是自出题自评测 & 1 天 \\
录一段完整演示 & 五条链路各有一段可播放记录 & 半天 \\
\bottomrule
\end{tabular}
\end{table}

最后一项需要说明：这台机器上没有可见桌面和音频，我生成的是逐步截图序列加文字说明（在 \path{output/demo/} 与 \path{docs/DEMO-WALKTHROUGH.md}），它不是录像。要让评审看到真实操作过程，需要人在正常桌面环境下录一次。

\section{复现}

\begin{verbatim}
npm start                    # 启动，默认 127.0.0.1:8765
npm test                     # 221 项自动测试
npm run test:smoke           # headless Chrome 真打一局
npm run build:knowledge      # 从引擎重新生成知识卡
npm run eval:retrieval       # 检索对照实验
npm run train:intervention   # 干预策略实验
\end{verbatim}

来源清单在 \code{knowledge/sources.json} 与 \code{docs/DOMAIN-RESEARCH.md} 两份文件里，官方文档、公开转载与我的推断分开标注，不混在一起。
